\chapter{Theory Background}

\section{Probability theory}

\subsection{Measure and integration}

\begin{definition}
    Let $\Omega$ be a set and let $2^\Omega$ denote its power set. A subset $\mathcal{F} \subset 2^\Omega$ is called a $\sigma$-\textbf{algebra
    over} $\Omega$ if it satisfies the following three conditions:
    \begin{enumerate}
        \item $\emptyset \in \mathcal{F}$.
        \item If $A, B \in \mathcal{F}$, then $B \backslash A \in \mathcal{F}$.
        \item For any countable sequence $A_1, A_2, \dots \in \mathcal{F}$, we have $\bigcup_{n=1}^\infty A_n \in \mathcal{F}$.
    \end{enumerate}
    If $\mathcal{F}$ is a $\sigma$-algebra over $\Omega$, then the tuple $(\Omega, \mathcal{F})$
    is called a \textbf{measurable space}. For any subset $\mathcal{G} \subset 2^\Omega$, we define
    the $\sigma$-\textbf{algebra generated by} $\mathcal{G}$ as the intersection over all $\sigma$-algebras
    that contain $\mathcal{G}$ as an element, and we denote this $\sigma$-algebra by $\sigma(\mathcal{G})$.
\end{definition}

\begin{example}
    \label{ex:borel-sigma-algebra}
    An important example of a $\sigma$-algebra over $\R^d$ is the \textbf{Borel} $\sigma$-\textbf{algebra}
    $\mathcal{B}(\R^d)$, which is defined to be the $\sigma$-algebra generated by the subset of all
    open sets on $\R^d$. Functions that are $\mathcal{B}(\R^d)$-measurable are called \textbf{Borel measurable}.
\end{example}

Unless otherwise stated, we will always consider $\R^d$ to be endowed with the Borel $\sigma$-algebra.

\begin{definition}
    Let $(\Omega, \mathcal{F})$ and $(E, \mathcal{G})$ be measurable spaces.
    A map $f\colon \Omega \to E$ is called $\mathcal{F},\mathcal{G}$-\textbf{measurable}
    if $f^{-1}(G) := \{ \omega\in\Omega \, | \, f(\omega) \in G \} \in \mathcal{F}$
    for all $G\in\mathcal{G}$.  We may say $f$ is $\mathcal{F}$-\textbf{measurable}
    or simply \textbf{measurable} if one or both of the $\sigma$-algebras are either
    clear from the context or not relevant.
\end{definition}

\begin{definition}
    Let $(\Omega, \mathcal{F})$ be a measurable space. A map $\mu : \mathcal{F} \to [0, \infty]$ is called a \textbf{measure on} $(\Omega, \mathcal{F})$ if it satisfies the following two conditions:
    \begin{enumerate}
        \item $\mu(\emptyset) = 0$.
        \item For any countable sequence $A_1, A_2, \dots \in \mathcal{F}$, we have $\mu(\bigcup_{n=1}^{\infty} A_n) = \sum_{n=1}^{\infty} \mu(A_n)$.
    \end{enumerate}
    If $\mu$ is a measure on $(\Omega, \mathcal{F})$, then the triplet $(\Omega, \mathcal{F}, \mu)$ is called a \textbf{measure space}.
\end{definition}

\begin{definition}
    Let $(\Omega, \mathcal{F}, \mu)$ be a measure space and let $f\colon \Omega \to \R$ be measurable. We define the ($\mu$-)\textbf{integral} of $f$, denoted $\int f \, \textnormal{d}\mu$, in three steps:
    \begin{enumerate}
        \item If $f(\omega) = \sum_{i=1}^n c_i 1_{A_i}(\omega)$ for some $n\in\N$, $c_1, \dots, c_n > 0$, and disjoint measurable sets $A_1, \dots, A_n\in\mathcal{F}$, then we define
        \begin{align*}
            \int f \, \textnormal{d}\mu := \sum_{i=1}^n c_i \mu(A_i).
        \end{align*}

        In this case $f$ is called a \textbf{simple function}. The set of all simple functions on $\Omega$ is denoted by $\mathcal{S}(\Omega)$.

        \item If $f$ is nonnegative, i.\,e. $f(\omega) \geq 0$ of all $\omega\in\Omega$, then we define
        \begin{align*}
            \int f \, \textnormal{d}\mu := \sup_{g \in \mathcal{S}(\Omega), \, g \leq f} \int g \, \textnormal{d}\mu.
        \end{align*}

        \item If $f$ is neither a simple function, nor nonnegative, but $\int |f| \, \textnormal{d}\mu < \infty$, then we define
        \begin{align*}
            \int f \, \textnormal{d}\mu := \int \max(0, f) \, \textnormal{d}\mu - \int \max(0, -f) \, \textnormal{d}\mu.
        \end{align*}
    \end{enumerate}
    Otherwise, we say that the ($\mu$-)integral of $f$ does not exist. If any of these three conditions apply to $f$, we say that $f$ is ($\mu$-)\textbf{integrable}.
\end{definition}

Apart from notation used above, we will also sometimes use
\[
    \int f(x) \, \mu(\text{d}x)
\]
to denote the $\mu$-integral of $f$.

\begin{proposition}
    \label{prop:integrals}
    Let $(\Omega, \mathcal{F}, \mu)$ be a measure space and let $f\colon \Omega \to \R$ and
    $g\colon \Omega \to \R$ be integrable. Then, the following hold:
    \begin{enumerate}
        \item[(i)] $\int f + g \, \textnormal{d}\mu = \int f \, \textnormal{d}\mu + \int g \, \textnormal{d}\mu$.
        \item[(ii)] $\int c f \, \textnormal{d}\mu = c \int f \, \textnormal{d}\mu$ for all $c\in\R$.
        \item[(iii)] If $f \leq g$, then $\int f \, \textnormal{d}\mu \leq \int g \, \textnormal{d}\mu$. If additionally $f < g$ on some set $A\in\mathcal{F}$ with $\mu(A) > 0$, then $\int f \, \textnormal{d}\mu < \int g \, \textnormal{d}\mu$.
    \end{enumerate}
\end{proposition}

\begin{proposition}
    \label{prop:integral-transform}
    Let $(\Omega, \mathcal{F}, \mu)$ be a measure space and let $f\colon \Omega \to \R$ and
    $g\colon \R \to \R$ be integrable. Then, it holds that
    \[
        \int g \circ f \, \text{d}\mu = \int g \, \text{d}\mu^{f},
    \]
    where $\mu^f$ is the measure defined as
    \[
        \mu^{f}(B) := \mu(f^{-1}(B)),
    \]
    for all $B\in \mathcal{B}(\R)$.
\end{proposition}

\subsection{Random variables, expected value, and variance}

\begin{definition}
    Let $(\Omega, \mathcal{F})$ be a measurable space. We say that a measure
    $\Prob: \mathcal{F} \to [0, \infty)$ is a \textbf{probability measure} measure on
    $(\Omega, \mathcal{F})$, if
    \begin{enumerate}
        \item $\Prob(A) \in [0, 1]$ for all $A\in\mathcal{F}$, and
        \item $\Prob(\Omega) = 1$.
    \end{enumerate}
    If $\Prob$ is a probability measure on $(\Omega, \mathcal{F})$, then we
    call the triplet $(\Omega, \mathcal{F}, \Prob)$ a \textbf{probability space}.
    In this context, elements of $\mathcal{F}$ are called \textbf{events}.
\end{definition}

\begin{definition}
    Let $(\Omega, \mathcal{F}, \Prob)$ be a probability space. An event $A\in\mathcal{F}$ is said to hold \textbf{almost surely} (a.\,s. for short) if $\Prob(A) = 1$.
\end{definition}

\begin{definition}
    Let $(\Omega, \mathcal{F}, \Prob)$ be a probability space and $(E, \mathcal{G})$ a measurable space. A map $X\colon \Omega \to E$ is called
    a \textbf{random variable} on $(\Omega, \mathcal{F}, \Prob)$ if $X$ is $\mathcal{F},\mathcal{G}$-measurable.
    In the case $E = \R^n$, we may call $X$ a \textbf{random vector}.
    Further, we define the notation $\Prob(X \in G) := \Prob(X^{-1}(G))$.
    We define the \textbf{distribution of} $X$ to be the probability measure
    $\Prob^X = \Prob \circ X^{-1}$ on $(E, \mathcal{G})$. Finally, we define
    $\sigma(X) := \sigma(\{ X^{-1}(G),\, G\in \mathcal{G} \})$
    and call this the \textbf{$\sigma$-algebra
    generated by $X$}.
\end{definition}

\begin{definition}
    Let $(\Omega, \mathcal{F}, \Prob)$ be a probability space. Two random variables $X$ and $Y$
    on $(\Omega, \mathcal{F}, \Prob)$ are called \textbf{independent} if $\Prob(X\in A,\, Y\in B)
    = \Prob(X\in A)\cdot \Prob(Y\in B)$ for all $A,B\in\mathcal{F}$. $X$ and $Y$ are called
    \textbf{identically distributed} if $\Prob^X = \Prob^Y$.
    We use the abbreviation \textbf{i.\,i.\,d.} as shorthand for "independent and identically distributed".
\end{definition}

\begin{definition}
    Let $(\Omega, \mathcal{F}, \Prob)$ be a probability space and let $X\colon\Omega\to\R$
    be a random variable on this space.
    If $X$ is integrable, we define the \textbf{expected value of} $X$,
    denoted by $\E(X)$, as
    \[
        \E(X) := \int X \, \textnormal{d}\Prob.
    \]
    If $X\colon\Omega\to\R^n$ is a random vector with components
    $X_i\colon\Omega\to\R$ for $i\in\{\, 1,\dots,n \,\}$, then we
    say that \textbf{$X$ is integrable} if all components are integrable. In that case,
    we define \textbf{the expected
    value of $X$} as the vector
    \[
        \E(X) := (\E(X_1),\dots, \E(X_n))^\top.
    \]
\end{definition}
The following three properties will be used multiple times throughout this text without explicit mention. They
follow directly from \cref{prop:integrals}.
\begin{proposition}
    \label{prop:properties-of-expected-value}
    Let $(\Omega, \mathcal{F}, \Prob)$ be a
    probability space and let $X\colon \Omega \to \R^n$ and
    $Y\colon \Omega \to \R^n$ be integrable random variables. Then, the following hold:
    \begin{enumerate}
        \item[(i)] $\E(X + Y) = \E(X) + \E(Y)$.
        \item[(ii)] $\E(c X) = c\, \E(X)$ for all $c\in\R$.
        \item[(iii)] If $X \leq Y$, then $\E(X) \leq \E(Y)$. If, additionally,
        $X(\omega) < Y(\omega)$ for all $\omega$ in an event $A\in\mathcal{F}$ with $\Prob(A) > 0$,
        then $\E(X) < \E(Y)$.
    \end{enumerate}
\end{proposition}

% \begin{proposition}
%     \label{prop:expectation-preserves-convexity}
%     Let $f\colon \R^d\times\Omega \to\R$ be convex in its first argument, i.\,e. $x~\mapsto~f(x, \omega)$
%     is convex for all $\omega\in\Omega$. Then, the function $x\mapsto \E(f(x, \cdot))$ is convex. 
% \end{proposition}

\begin{proposition}[Jensen's inequality]
    \label{prop:jensens-inequality} 
    Let $(\Omega, \mathcal{F}, \Prob)$ be a probability space and let $X\colon \Omega\to\R^n$
    be a random variable. Then we have
    \[
        \E(X^2) \geq \E(X)^{2}.
    \]
    In particular: If $X^2$ is integrable, then $X$ must also be integrable.
\end{proposition}

\begin{definition}
    Let $(\Omega, \mathcal{F}, \Prob)$ be a probability space and let $X\colon \Omega \to \R^n$ be a random variable.
    If $X^2$ is integrable, we define the \textbf{variance of} $X$, denoted by $\Var(X)$, as $\Var(X) := \E|\!|X-\E(X)|\!|^2$.
\end{definition}
By \cref{prop:integral-transform}, for any integrable random variable 
$X: \Omega\to\R$, it holds that $\int X \,\textnormal{d}\Prob = \int I \,\textnormal{d}\Prob^X$, where $I\colon \R\to \R$ is the identity
operator. It is now easy to see that if two random variables $X$ and $Y$ are identically distributed,
then they have the same expected value and variance.
\begin{proposition}
    \label{prop:variance-linear-for-independent-rvs}
    Let $(\Omega, \mathcal{F}, \Prob)$ be a probability space and let $X_1,\dots,X_N\colon \Omega \to \R^n$
    be $N\in\N$ independent random variables, such that $\norm{X_i}^2$ is integrable
    for all $i\in\{1,\dots, N\}$. Then $\Var(\sum_{i=1}^N a_i X_i) = \sum_{i=1}^{N} a_i^2\, \Var(X_i)$,
    for any $a_1,\dots, a_N\in\R$.
    If, additionally, they are all identically distributed, it holds that
    $\Var(1/N \sum_{i=1}^N X_i) = \Var(X_1)/N$.
\end{proposition}

\begin{proposition}
    \label{prop:variance-identity}
    Let $(\Omega, \mathcal{F}, \Prob)$ be a probability space and let $X\colon \Omega \to \R^n$
    be a random variable such that $X^2$ is integrable. Then, $\Var(X) = \E(X^2) - \E(X)^2$.
\end{proposition}

\subsection{Densities}

\begin{definition}[Density]
    Let $(E, \mathcal{E})$ be a measurable space, and let $\mu$ and $\nu$
    be measures on $(E, \mathcal{E})$.
    A function $f\colon E\to [0,\infty)$
    is called a \textbf{$\nu$-density of $\mu$}, if
    \[
        \mu(A) = \int_{A} f \, \textnormal{d}\nu,
    \]
    for all $A\in \mathcal{E}$. In the special case where $E = \R^n$ and
    $\nu$ is the Lebesgue-measure on $\R^n$,
    we refer to $f$ as a \textbf{Lebesgue-density}, or simply \textbf{density}, of $\mu$.
\end{definition}

\begin{definition}[Absolute continuity and equivalence]
    Let $(\Omega, \mathcal{F})$ be a measurable space and let $\mu$ and $\nu$ be measures
    on $(\Omega, \mathcal{F})$.
    We say that $\mu$ is \textbf{absolutely continuous with respect to $\nu$},
    if
    \[
        \nu(A) = 0 \quad\text{implies}\quad \mu(A) = 0,
    \]
    for all $A\in\mathcal{F}$. If $\mu$ is absolutely continuous with respect to $\nu$,
    then we use the notation $\mu \ll \nu$.
    We say that $\mu$ and $\nu$
    are \textbf{equivalent}, if $\mu \ll \nu$ and $\nu \ll \mu$.
    If $\mu$ and $\nu$ are equivalent, then we use the notation $\mu \sim \nu$.
\end{definition}

\begin{proposition}[Radon-Nikodym]
    \label{prop:radon-nikodym}
    Let $(\Omega, \mathcal{F})$ be a measurable space and let $\mu$ and $\nu$ be measures
    on $(\Omega, \mathcal{F})$. Then, $\mu \ll \nu$
    if and only if there exists a $\nu$-density of $\mu$.
    If $\mu$ and $\nu$ are equivalent and $f$ is a $\nu$-density of $\mu$, then $f > 0$
    $\nu$-almost everywhere
    and there exists a $\mu$-density $g$ of $\nu$, such that $f \cdot g = 1$ $\nu$-almost everywhere.
\end{proposition}
\begin{proposition}
    \label{prop:density-chain-rule}
    Let $(E, \mathcal{E})$ be a measurable space and let $\mu$, $\nu$ and $\tau$
    be measures on $(E, \mathcal{E})$. Suppose that $\mu$ has a $\nu$-density $f$ and $\nu$ has a $\tau$-density
    $g$. Then, the function $h := f \cdot g$ is a $\tau$-density of $\mu$.
\end{proposition}
If $f$ is a $\nu$-density of $\mu$, then $f$ is often referred to as a \textbf{Radon-Nikodym derivative of $\mu$
with respect to $\nu$}. Commonly used notations are
\[
    f = \frac{\text{d}\mu}{\text{d}\nu}
\]
and
\[
    \text{d}\mu = f \text{d}\nu.
\]
Putting the two together, we formally have
\[
    \text{d}\mu = \frac{\text{d}\mu}{\text{d}\nu} \text{d}\nu.
\]
This notation suggest that, if $\mu \ll \nu$ and $\nu \ll \tau$ for some other measure $\tau$, then we have
\[
    \text{d}\mu = \frac{\text{d}\mu}{\text{d}\nu} \text{d}\nu
    = \frac{\text{d}\mu}{\text{d}\nu} \frac{\text{d}\nu}{\text{d}\tau} \text{d}\tau,
\]
and hence $\mu \ll \tau$ with density $\frac{\text{d}\mu}{\text{d}\nu} \frac{\text{d}\nu}{\text{d}\tau}$,
which is consistent with the above proposition.

We now move to back to the setting of probability spaces, where the concept of
densities are of great help to work with random variables.
\begin{definition}[Density of a random variable]
    Let $(\Omega, \mathcal{F}, \Prob)$ be a probability space and let $X\colon \Omega \to \R^n$ be a
    random variable. Further, let $\mu$ be a measure on $(\Omega, \mathcal{F})$.
    A function $f\colon\R\to [0,\infty)$ is called a \textbf{$\mu$-density of $X$},
    if $f$ is a $\mu$-density of $\Prob^X$. That is,
    \[
        \Prob(X \in A) = \int_{A} f \, \textnormal{d}\mu,
    \]
    for all $A\in\mathcal{B}(\R^n)$.
\end{definition}

The following is an important special case of \cref{prop:density-chain-rule}.
\begin{proposition}
    Let $(\Omega, \mathcal{F}, \Prob)$ be a probability space and let $X\colon \Omega \to \R^n$ be a
    random variable. Further, let $\mu$ be a measure on $(\Omega, \mathcal{F})$ and suppose
    that $X$ has a $\mu$-density $f$ such that $\int |x|f(x) \, \text{d}x < \infty$. Then, $X$ is integrable
    and
    \[
        \E(X) = \int x f(x) \, \mu(dx).
    \]
    In particular, if $\mu$ is the Lebesgue-measure on $\R^n$, it holds that
    \[
        \E(X) = \int x f(x) \, \textnormal{d}x.
    \]
    If $\mu$ is instead the counting measure on a discrete subset $S\subset \R^n$, then
    \[
        f(x) = \Prob(X = x),
    \]
    for all $x\in S$, and
    \[
        \E(X) = \sum_{x\in S} x \, \Prob(X = x).
    \]
\end{proposition}

\subsection{Conditioning}

\begin{definition}
    Let $(\Omega, \mathcal{F}, \Prob)$ be a probability space and let $X\colon \Omega\to\R^n$ be a random variable.
    Further, let $\mathcal{C}\subset\mathcal{F}$ be a $\sigma$-algebra.
    We call a random variable $Z\colon \Omega\to\R^n$ a \textbf{conditional expectation of} $X$ \textbf{given} $\mathcal{C}$, if
    \begin{enumerate}
        % \item There exists some measurable $h\colon \R^m\to\R^n$ such that $Z = h(Y)$.
        \item $Z$ is $\mathcal{C}$-measurable, and
        \item for all $C\in \mathcal{C}$, it holds that
        \begin{equation*}
            \int_{C} Z \, \textnormal{d}\Prob = \int_{C} X \, \textnormal{d}\Prob.
        \end{equation*}
    \end{enumerate}
    If $Z$ is a conditional expectation of $X$ given $\mathcal{C}$ then we use the notation $\E(X \mid \mathcal{C}) := Z$.
    
    If $Y\colon\Omega\to \R^m$ is a random variable such that $\mathcal{C} = \sigma(Y)$, then
    we use the notation $\E(X \mid Y) := \E(X \mid \sigma(Y))$.
    In that case, we call $\E(X \mid Y)$ the \textbf{conditional expectation of $X$ given $Y$}.
    Further, for $\omega\in\Omega$ with $Y(\omega) = y\in\R^m$, the \textbf{conditional expectation of $X$ given $Y=y$}, denoted
    by $\E(X \mid Y = y)$, is defined as $\E(X \mid Y = y) := \E(X \mid Y)(\omega)$.
\end{definition}
Note that $E(X \mid Y)$ is not unique. However, if $Z_1$ and $Z_2$ are both conditional expectations
of $X$ given $Y$, then we always have $Z_1 = Z_2$ almost surely. For simplicity, we will keep the
"almost surely" implicit.
\begin{remark}
    \label{remark:conditional-expectation}
    If $X$ and $Y$ are integrable random variables on a probability space $(\Omega, \mathcal{F}, \Prob)$, then
    the conditional expectation of $X$ given $Y=y$ can be thought of as the expected value of
    $X$ on a different probabiliy space $(\Omega, \mathcal{F}, \Prob_{Y=y})$, where
    $\Prob_{Y=y}(A) := \Prob(A \mid Y=y)$. More precisely,
    $\E(X \mid Y = y) = \int X \, \textnormal{d}\Prob_{Y=y}$. It follows that $\E(X \mid Y = \cdot)$
    inherits all basic properties of $\E(\cdot)$ from \cref{prop:properties-of-expected-value}.
\end{remark}
Below, we state some further properties of conditional expectations.

\begin{proposition}[Properties of $\E(X \mid Y)$]
    \label{prop:conditional-expectation-properties}
    Let $(\Omega, \mathcal{F}, \Prob)$ be a probability space and
    let $X\colon \Omega \to \R^n$, $Y\colon \Omega \to \R^m$ be integrable random variables. Further, let $F\colon\R^n\times \Omega\to\R$ such that
    $\omega\mapsto F(x, \omega)$ is $\mathcal{F}$-measurable for all $x\in\R^d$ and let $G\colon\R^n\to\R$
    be Borel measurable. Then, the following hold:
    \begin{itemize}
        \item[\textnormal{(i)}] $\E(\E(X \mid Y)) = \E(X)$.
        \item[\textnormal{(ii)}] $\E(X \mid X) = X$.
        \item[\textnormal{(iii)}] If $F(x, \cdot) \leq G(x)$ a.\,s. for all $x\in\R^d$, it follows that $\E(F(X, \cdot) \mid X) \leq G(X)$.
    \end{itemize}
\end{proposition}

\begin{definition}
    Let $(\Omega, \mathcal{F}, \Prob)$ be a probability space and
    let $X\colon \Omega \to \R^n$, $Y\colon \Omega \to \R^m$ be integrable random variables.
    The \textbf{conditional variance of X given Y}, denoted by $\Var(X \mid Y)$, is defined as
    \[
        \Var(X \mid Y) := \E\big( ||X - \E(X \, | \, Y)||^2 \mid Y \big).
    \]
    For $y\in\R^d$, the \textbf{conditional variance of X given Y = y} is defined as
    \[
        \Var(X \mid Y = y) := \E\big( ||X - \E(X \, | \, Y = y)||^2 \mid Y = y \big).
    \]
\end{definition}

\begin{remark}
    Let $X$ and $Y$ be integrable random variables on a probability space $(\Omega, \mathcal{F}, \Prob)$.
    Similarly to the conditional expectation of $X$ given $Y = y$, the conditional variance
    $\Var(X \mid Y = y)$ can be thought of as the variance of $X$ on a different probability
    space $(\Omega, \mathcal{F}, \Prob_{Y=y})$ (see \cref{remark:conditional-expectation}).
    Hence, all basic properties of $\Var(\cdot)$ are inherited by the conditional variance.
    In particular, if $X_1,\dots, X_N$ are i.\,i.\,d. random variables, it holds that
    $\Var(X_1 + \cdots + X_N \mid Y) = \Var(X_1 \mid Y) / N$ -- a fact that will be used later.
\end{remark}

\begin{proposition}
    \label{prop:conditional-variance-properties}
    Let $(\Omega, \mathcal{F}, \Prob)$ be a probability space and
    let $X\colon \Omega \to \R^n$ be an integrable random variable.
    Further, let $F:\R^n\times \Omega\to\R$ such that
    $\omega\mapsto F(x, \omega)$ is $\mathcal{F}$-measurable for all $x\in\R^d$ and $\E(F(X,\cdot)^2) < \infty$, and let $G:\R^n\to\R$
    be Borel measurable. Then, if $\Var(F(x, \cdot)) \leq G(x)$ for all $x\in\R^d$, it also holds that
    \[
        \Var(F(X, \cdot) \mid X) \leq G(X).
    \]
\end{proposition}

\subsection{Martingales}

\begin{definition}[Filtration]
    Let $(\Omega, \mathcal{F}, \Prob)$ be a probability space. A sequence
    of $\sigma$-algebras $(\mathcal{F}_n)_{n\in\N_0}$ with
    $\mathcal{F}_n \subset \mathcal{F} \, \forall n\in\N_0$
    is called a \textbf{filtration} if $\mathcal{F}_{n-1} \subset \mathcal{F}_{n}$
    for all $n\in\N$. If $(\mathcal{F}_n)_{n\in\N_0}$ is a filtration on
    $(\Omega, \mathcal{F}, \Prob)$, then we call
    $(\Omega, \mathcal{F}, \Prob, (\mathcal{F}_n)_{n\in\N_0})$ a \textbf{filtered probability space}.
\end{definition}

\begin{example}[Natural filtration]
    Given a sequence of random variables $(X_n)_{n\in\N_0}$ on some common probabilty
    space $(\Omega, \mathcal{F}, \Prob)$, one can define a filtration
    $(\mathcal{F}_n)_{n\in\N}$ via $\mathcal{F}_n := \sigma(X_0, \dots, X_n)$ for all $n\in\N_0$.
    This filtration is referred to as the \textbf{natural filtration}
    of the sequence $(X_n)_{n\in\N_0}$.
\end{example}

\begin{definition}
    Let $(\Omega, \mathcal{F}, \Prob)$ be a probability space. A
    sequence of random variables $(X_t)_{t\in T}$, where $T \subset \R$
    is some set of indeces, is called a
    \textbf{stochastic process}.
    Let $(\mathcal{F}_n)_{n\in\N_0}$ be a filtration on
    $(\Omega, \mathcal{F}, \Prob)$. A stochastic process $X := (X_n)_{n\in\N_0}$
    is called \textbf{adapted to $(\mathcal{F}_n)_{n\in\N_0}$}, if
    $X_n$ is $\mathcal{F}_n$-measurable for all $n\in\N_0$. 
\end{definition}

Clearly, any stochastic process is adapted to their natural filtration.

\begin{definition}[Martingale]
    Let $X := (X_n)_{n\in\N}$ be a stochastic process on some
    common filtered probability space
    $(\Omega, \mathcal{F}, \Prob, (\mathcal{F}_n)_{n\in\N_0})$.
    Then, we say that $X$ is a \textbf{martingale}, if the following conditions are satisfied:
    \begin{enumerate}
        \item $\E|X_n| < \infty$ for all $n\in\N_0$.
        \item $X$ is adapted to $(\mathcal{F}_n)_{n\in\N_0}$, i.\,e. $X_n$ is $\mathcal{F}_n$-measurable for all $n\in\N_0$.
        \item $\E(X_{n+1} \mid X_{n}) = X_n$ for all $n\in\N_0$.
    \end{enumerate}
\end{definition}

\begin{proposition}[Azuma-Hoeffding inequality]
    \label{prop:azuma-hoeffding}
    Let $(\Omega, \mathcal{F}, \Prob, (\mathcal{F}_n)_{n\in\N_0})$
    be a filtered probability space
    and let $(X_k)_{k\in\N_0}$ be a martingale and
    let $(c_k)_{k\in\N}$ be a sequence such that
    \[
        |X_k - X_{k-1}| \leq c_k
    \]
    almost surely, for all $k\in\N$. Then, for all $K\in\N$
    and all $t\in (0, \infty)$, it holds that
    \[
        \Prob(X_K - X_0 \geq t) \leq \exp\Big( \frac{-t^2}{2 \sum_{k=1}^K c_k^2} \Big).
    \]
\end{proposition}

\section{Optimization theory}
\subsection{Lipschitz functions}
\subsection{Convex optimization}
\subsection{Stochastic optimization}
\subsection{Multifunctions and metric regularity}

\begin{definition}[Multifunctions]
    Let $X$ and $Y$ be Banach spaces. A function $\varPsi\colon X\to 2^Y$
    is called \textbf{multifunction}. The \textbf{domain} and the \textbf{range}
    of a multifunction $\varPsi\colon X\to 2^Y$ are defined as
    \begin{align*}
        \dom(\varPsi) &:= \{\, x\in X \mid \varPsi(x) \neq \emptyset \,\}, \\
        \range(\varPsi) &:= \{\, y\in Y \mid y\in \varPsi(x) \text{ for some $x\in X$} \,\}.
    \end{align*}
    The \textbf{graph} of $\varPsi$ is
    \[
        \graph(\varPsi) := \{\, (x, y) \in X\times Y \mid y\in\varPsi(x), x\in X \,\}.
    \]
    The \textbf{(graph) inverse} $\varPsi^{-1}\colon Y \to 2^X$ of $\varPsi$ is defined as
    \[
        \varPsi^{-1}(y) := \{\, x\in X \mid y\in\varPsi(x) \,\}.
    \]
    The multifunction $\varPsi$ is said to be \textbf{closed at a point} $x\in X$ if
    $x_k \to x$, $y_k\in \varPsi(x_k)$, and $y_k \to y$ imply that $y\in\varPsi(x)$.
    If $\varPsi$ is closed at all points $x\in X$, then $\varPsi$ is said to be \textbf{closed}.
    Finally, it is said that $\varPsi$ is \textbf{convex} if, for any $x_1,x_2\in X$ and
    $t\in [0, 1]$,
    \[
        t \varPsi(x_1) + (1 - t) \varPsi(x_2) \subset \varPsi(t x_1 + (1 - t) x_2).
    \]
\end{definition}
\begin{definition}[Metric regularity]
    We say that the multifunction $\varPsi\colon X\to 2^Y$ is \textbf{metric regular} at a point
    $(x_0,y_0)\in\graph(\varPsi)$ at rate $c\in\R_{\geq 0}$, if there
    exists a neighborhood $U\subset X\times Y$ containing $(x_0,y_0)$, such that
    \[
        \dist(x, \varPsi^{-1}(y)) \leq c \dist(y, \varPsi(x)),
    \]
    for all $(x, y)\in U$.
\end{definition}
\begin{proposition}[Robinson-Ursescu stability theorem]
    \label{prop:robinson-ursescu}
    Let $\varPsi\colon X\to 2^Y$ be a closed convex multifunction. Then $\varPsi$
    is metric regular at $(x_0, y_0) \in \graph(\varPsi)$ if and only if the regularity
    condition $y_0 \in \interior (\range \varPsi)$ holds.
\end{proposition}
